\documentclass[11pt]{article}
\usepackage[margin=1in]{geometry}
\usepackage{amsmath,amssymb,amsthm}
\usepackage{graphicx}
\usepackage[hidelinks]{hyperref}
\usepackage{booktabs}

\newtheorem{theorem}{Theorem}
\newtheorem{lemma}{Lemma}
\newtheorem{corollary}{Corollary}
\newcommand{\N}{\mathbb N}
\newcommand{\Id}{\operatorname{Id}}
\newcommand{\cspan}{\overline{\operatorname{span}}}
\newcommand{\ellone}{\ell^{1}}
\newcommand{\ellinf}{\ell^{\infty}}

\title{An Unbounded Diagonal Semigroup on the Bourgain--Delbaen Space\\
and the Dual of Its Sun Dual}
\author{Automatic Math Discovery Run \texttt{fa\_banach\_001}}
\date{July 19, 2026}

\begin{document}
\maketitle

\begin{center}
\fbox{\parbox{0.92\textwidth}{\textbf{Status.} Candidate full solution,
likely valid, subject to expert review.  The construction has an unbounded
generator and therefore meets the strongest natural interpretation of
``non-trivial'' in the source question.}}
\end{center}

\section*{Source and Extracted Question}

The source is Philip Preussler and Felix L. Schwenninger,
\emph{Implications of structured continuous maximal regularity},
arXiv:2605.12121v1 \cite{preussler-schwenninger}.  In Remark 4.15(2), page 24,
the authors consider the classical Bourgain--Delbaen space \(Y\), which
satisfies
\[
  Y\not\supset c_0,
  \qquad Y'\simeq \ellone,
  \qquad Y''\simeq \ellinf\supset c_0,
\]
and ask whether there exists a non-trivial \(C_0\)-semigroup on \(Y\) such
that
\[
  (Y^\odot)'\supset c_0.
\]
Here \(Y^\odot\) is the sun dual associated with the semigroup.  The source
immediately discusses the possibility that \(Y\) has the Lotz property, i.e.
that every \(C_0\)-semigroup on \(Y\) has bounded generator.  Thus the strongest
and mathematically substantive reading of ``non-trivial'' is ``having
unbounded generator.''  The construction below has that property.

\begin{center}
\includegraphics[width=\textwidth]{figures/open_problem_crop.png}
\end{center}

\section*{Definitions}

Let \(X\) be a Banach space and \(T=(T(t))_{t\geq0}\) a \(C_0\)-semigroup on
\(X\).  Its \emph{sun dual} is
\[
 X^\odot
 =\left\{x'\in X':\lim_{t\downarrow0}\|T(t)'x'-x'\|=0\right\}.
\]
It is a norm-closed \(T(t)'\)-invariant subspace of \(X'\).

If \((x_n)\) is a Schauder basis, write \(x_n'\) for its biorthogonal
functionals and
\[
 P_Nx=\sum_{n=1}^N x_n'(x)x_n,
 \qquad K=\sup_N\|P_N\|<\infty.
\]

\section*{Proof Intuition}

There are two independent mechanisms.  First, any Schauder basis supports a
bounded diagonal heat-type semigroup: multiply its \(n\)-th basis vector by
\(e^{-nt}\).  Monotonicity of these scalars lets Abel summation control the
operator norm even when the basis is conditional.  The eigenvalues \(-n\)
force the generator to be unbounded.

Second, every coordinate functional is an eigenvector of the adjoint
semigroup and hence belongs to the sun dual.  Thus the sun dual is an
infinite-dimensional subspace of \(Y'\simeq\ellone\).  A gliding-hump
argument shows that every such subspace contains a complemented copy of
\(\ellone\).  Dualizing that complemented summand puts a complemented copy of
\(\ellinf\), and therefore \(c_0\), inside \((Y^\odot)'\).

The proof never requires the Bourgain--Delbaen basis to be unconditional or
shrinking.

\section*{A Complemented \(\ellone\) Lemma}

We include the classical structural fact needed in the proof so that the only
external dependency is the identification of the Bourgain--Delbaen space.

\begin{lemma}[Complemented \(\ellone\)-saturation]
\label{lem:l1}
Every infinite-dimensional closed subspace \(E\) of \(\ellone\) contains a
closed subspace \(W\simeq\ellone\) that is complemented in \(E\).  In fact,
\(W\) may be chosen complemented in the ambient \(\ellone\).
\end{lemma}

\begin{proof}
Let \(R_N:\ellone\to\ellone\) be truncation to the first \(N\) coordinates.
Choose positive numbers \((\varepsilon_k)\) with
\(\delta:=2\sum_k\varepsilon_k<1\).  Inductively choose integers
\(0=N_0<N_1<N_2<\cdots\) and unit vectors \(z_k\in E\) such that
\[
 R_{N_{k-1}}z_k=0,
 \qquad
 \|(\Id-R_{N_k})z_k\|_1<\varepsilon_k.
\]
This is possible because \(R_{N_{k-1}}|_E\) has finite rank, so its kernel in
\(E\) is infinite-dimensional, and every vector of \(\ellone\) has small
tails.

Put
\[
 u_k=(R_{N_k}-R_{N_{k-1}})z_k,
 \qquad v_k=\frac{u_k}{\|u_k\|_1}.
\]
The vectors \(v_k\) have pairwise disjoint supports and norm one, and
\[
 \|z_k-v_k\|_1
 \leq \|z_k-u_k\|_1+|1-\|u_k\|_1|
 <2\varepsilon_k.
\]
Define \(U:\ellone\to\ellone\) by
\(U(a)=\sum_k a_kv_k\).  Disjointness gives \(\|U(a)\|_1=\|a\|_1\).
For every \(k\), choose a norm-one functional \(\phi_k\) supported on the
support of \(v_k\) and satisfying \(\phi_k(v_k)=1\).  Then
\[
 P x=\sum_k\phi_k(x)v_k
\]
defines a contractive projection of \(\ellone\) onto
\(V:=U(\ellone)\).

Define \(S:\ellone\to\ellone\) by \(S(a)=\sum_k a_kz_k\).  The preceding
estimate gives
\[
 \|S-U\|\leq\sum_k\|z_k-v_k\|_1<\delta<1.
\]
Consequently \(S\) is an isomorphic embedding.  Moreover
\[
 B:=U^{-1}PS:\ellone\to\ellone
\]
satisfies \(\|B-\Id\|<1\), and hence is invertible.  The operator
\[
 Q:=SB^{-1}U^{-1}P:\ellone\longrightarrow S(\ellone)
\]
is bounded and satisfies \(Q(Sa)=Sa\) for every \(a\in\ellone\).  It is
therefore a projection onto
\(W:=S(\ellone)=\cspan\{z_k:k\in\N\}\).  Since every \(z_k\) lies in \(E\),
we have \(W\subset E\); restricting \(Q\) to \(E\) proves complementability in
\(E\) as well.
\end{proof}

\section*{The General Diagonal-Semigroup Theorem}

\begin{theorem}
\label{thm:general}
Let \(X\) be an infinite-dimensional Banach space with a Schauder basis and
suppose \(X'\simeq\ellone\).  Then \(X\) admits a \(C_0\)-semigroup \(T\) with
unbounded generator such that
\[
  (X^\odot)'\text{ contains a complemented copy of }\ellinf.
\]
In particular, \((X^\odot)'\) contains \(c_0\).
\end{theorem}

\begin{proof}
Let \((x_n)\) be a Schauder basis, with biorthogonal functionals
\((x_n')\) and basis projections \(P_N\).  For \(t\geq0\), define on basis
vectors
\[
  T(t)x_n=e^{-nt}x_n.
\]
For \(t>0\), Abel summation gives the operator-norm convergent formula
\[
  T(t)=\sum_{n=1}^{\infty}
       \bigl(e^{-nt}-e^{-(n+1)t}\bigr)P_n.
\]
Indeed, the \(N\)-th finite Abel sum has the additional remainder
\(e^{-Nt}P_N\), whose norm tends to zero.  Since all displayed coefficients
are nonnegative,
\[
 \|T(t)\|
 \leq K\sum_{n=1}^{\infty}
          \bigl(e^{-nt}-e^{-(n+1)t}\bigr)
 =Ke^{-t}.
\]
Set \(T(0)=\Id\).  The semigroup law holds on every \(x_n\), hence on the
dense finite-support span and then on all of \(X\).  On that finite-support
span, \(T(t)x\to x\) as \(t\downarrow0\).  The uniform bound
\(\sup_{0<t\leq1}\|T(t)\|\leq K\), followed by approximation with \(P_Nx\),
proves strong continuity at zero.  Thus \(T\) is a \(C_0\)-semigroup.

Let \(A\) be its generator.  For every \(n\),
\[
 \lim_{t\downarrow0}\frac{T(t)x_n-x_n}{t}=-n x_n,
\]
so \(x_n\in D(A)\) and \(Ax_n=-n x_n\).  Hence \(A\) is unbounded.

For each coordinate functional,
\[
 T(t)'x_n'=e^{-nt}x_n',
\]
and therefore \(x_n'\in X^\odot\).  Since the \(x_n'\) are linearly
independent, \(X^\odot\) is infinite-dimensional.  It is a closed subspace of
\(X'\simeq\ellone\), so Lemma~\ref{lem:l1} supplies a complemented subspace
\(W\subset X^\odot\) with \(W\simeq\ellone\).  Write
\(X^\odot=W\oplus Z\).  Dualizing the direct sum yields
\[
 (X^\odot)'\simeq W'\oplus Z',
 \qquad W'\simeq\ellinf.
\]
Thus \((X^\odot)'\) contains a complemented copy of \(\ellinf\), which in
turn contains its canonical copy of \(c_0\).
\end{proof}

\section*{Application to the Bourgain--Delbaen Question}

Tarbard's detailed exposition of the original Bourgain--Delbaen construction
shows that its biorthogonal vectors form a normalized Schauder basis
\cite[Chapter 2, especially the discussion following Proposition 2.4.3]
{tarbard}.  The same source records that every original Bourgain--Delbaen
space in the class \(\mathcal Y\), including the space denoted \(Y\) in the
question, has dual isomorphic to \(\ellone\)
\cite[Corollary 2.5.7]{tarbard}.

\begin{corollary}
The Bourgain--Delbaen space \(Y\) in Remark 4.15(2) of
arXiv:2605.12121 admits a \(C_0\)-semigroup with unbounded generator for which
\[
  (Y^\odot)'\supset \ellinf\supset c_0,
\]
where the copy of \(\ellinf\) is complemented.  Hence the question in that
remark has an affirmative answer.
\end{corollary}

\begin{proof}
Apply Theorem~\ref{thm:general} to the normalized Bourgain--Delbaen basis.
Explicitly, if that basis is \((y_n)\), use
\[
  T(t)y_n=e^{-nt}y_n.
\]
\end{proof}

\section*{Verification Notes}

\begin{enumerate}
\item \textbf{Conditionality causes no problem.}  The multiplier bound uses
only \(K=\sup_N\|P_N\|<\infty\) and the positive variation of
\((e^{-nt})_n\).  No unconditionality hypothesis is present.
\item \textbf{Shrinking is not used.}  The proof needs only the individual
coordinate functionals \(x_n'\) in the sun dual, not that they span all of
\(X'\).
\item \textbf{The generator is genuinely unbounded.}  The eigenvalue equation
\(Ax_n=-nx_n\) rules out a bounded generator directly.
\item \textbf{Isomorphic, not isometric, duality is enough.}  An isomorphism
\(X'\to\ellone\) sends \(X^\odot\) to an infinite-dimensional closed subspace;
complementability and the isomorphism type of the extracted \(\ellone\) copy
pull back.
\item \textbf{The conclusion is stronger than requested.}  The dual of the
sun dual contains complemented \(\ellinf\), rather than merely a copy of
\(c_0\).
\end{enumerate}

\section*{Novelty Search and Limitations}

On July 19, 2026, the lightweight indexes for run
\texttt{fa\_banach\_001} were searched for arXiv:2605.12121, its title,
maximal-regularity terms, input-to-state-stability terms, and the combinations
``Bourgain--Delbaen'', ``sun dual'', and ``Lotz property''.  No matching prior
solution or attempt was present.  Bounded web/arXiv searches used the exact
source title and combinations of ``Bourgain--Delbaen space'',
``\(C_0\)-semigroup'', ``sun dual'', ``unbounded generator'', and ``Lotz
property''.  They found the source and background construction papers but no
later explicit answer.  The arXiv record was still version 1, submitted May
12, 2026.

This search is not a global novelty certificate, and indexing can lag.  The
main review risk is bibliographic: an older general observation about diagonal
semigroups on \(\ellone\)-preduals may imply the same result without having
been phrased as an answer to this recent question.  The mathematical proof
above is independent of that novelty issue.

The result answers the Bourgain--Delbaen residual question only.  It does not
settle the separate general equivalence problem for
\(L^\infty\)-ISS and \(L^\infty\)-iISS in Remark 4.7 of the source.

\section*{Human-Review Recommendation}

The candidate should be sent to a Banach-space or operator-semigroup expert.
The most useful checks are: (i) matching the source's \(Y\) with Tarbard's
class \(\mathcal Y\); (ii) the gliding-hump projection in
Lemma~\ref{lem:l1}; and (iii) confirmation that ``non-trivial'' is intended in
the unbounded-generator sense suggested by the Lotz-property discussion.

\begin{thebibliography}{9}

\bibitem{preussler-schwenninger}
P.~Preussler and F.~L.~Schwenninger,
\emph{Implications of structured continuous maximal regularity},
arXiv:2605.12121v1, 2026.

\bibitem{tarbard}
M.~Tarbard,
\emph{Operators on Banach Spaces of Bourgain--Delbaen Type},
arXiv:1309.7469, 2013.

\bibitem{bourgain-delbaen}
J.~Bourgain and F.~Delbaen,
\emph{A class of special \(\mathcal L_\infty\) spaces},
Acta Math. 145 (1980), 155--176.

\end{thebibliography}

\end{document}
